\documentclass[11pt,a4paper]{article}

\usepackage[T1]{fontenc}
\usepackage[utf8]{inputenc}
\usepackage{amsmath,amssymb,amsthm}
\usepackage{geometry}
\geometry{margin=2.3cm}
\usepackage{booktabs}
\usepackage{graphicx}
\usepackage{hyperref}

\newcommand{\statut}[1]{{\footnotesize\textsf{[#1]}}}
\newcommand{\Mpl}{M_{\mathrm{Pl}}}
\newcommand{\Heff}{\hat{H}_{\mathrm{eff}}}
\newcommand{\Hsys}{\hat{H}_{\mathrm{sys}}}

\theoremstyle{definition}
\newtheorem{postulate}{Postulate}

\renewcommand{\thesection}{\Roman{section}}
\renewcommand{\thesubsection}{\thesection.\Alph{subsection}}

\title{\vspace{-1.0cm}\textbf{Unit 3 --- Compactified Scale Duality:\\
a geometric route toward the unification of quantum theory and gravitation}\\[5pt]
\large One scale field, from black holes to inflation, and a bridge to intermittent quantum existence}
\author{Hassane Bakkaoui\\ \small Independent researcher --- bakkahassa@hotmail.com}
\date{\small June 2026 $\cdot$ synthesis (gr-qc / astro-ph.CO / quant-ph)}

\begin{document}
\maketitle

\begin{abstract}
\noindent
The unification of quantum mechanics and general relativity remains open. We propose that
their two faces emerge from a single structure: the compactification of the scale axes,
spatial and temporal, into circles --- $\mathcal{M}=M_{3+1}\times S^1_\psi\times S^1_\chi$
under $\mathbb{Z}_2\times\mathbb{Z}_2$. A single scale field, coupled non-minimally via
$f=1+\beta(\cos\psi+\cos\chi)$ and carrying a bounded potential
$V=\alpha(2-\cos\psi-\cos\chi)$, governs both regimes. We establish: (i) natural inflation
at locked coupling, predictive --- $(n_s,r)$ and a constraint on $\beta$; (ii) the absence of
any regular black hole in the \emph{static} sector (closure), robust to un-freezing the
temporal axis; (iii) in collapse, that the verdict on whether curvature is capped is
\emph{reduction-dependent} and open --- it hinges on the nature of the pole $\infty^*$
(\S\ref{sec:collapse}); (iv) a bridge linking the temporal axis $\chi$ to the quantum time of
the Quantum Existence Discontinuity hypothesis (HDEQ), gravity modulating time.
The field-space metric is positive-definite ($f>0$, B1), the viable band the open interval
$\beta\in(-0.5;-0.15)$; the scale axes are internal circles ($\sigma$-model), with no
time-like dimension or ghost (F2); free propagation does not violate Lorentz. The existence
factor is derived exactly, $\varepsilon=E(k)/K(k)$. $\alpha$ is fixed by the scalar amplitude
$A_s$ ($\alpha\approx0.5$--$1.2\times10^{-9}\,\Mpl^4$, GUT scale), whence
$\tau_0\approx3.5\times10^{-37}$~s \emph{at the inflationary $f_s\approx5$--$6$} and a running
prediction $dn_s/d\ln k\approx-6.5\times10^{-4}$ (A1); the rescaling $\Heff=\varepsilon\Hsys$
is shown via a Page--Wootters relational clock (D1); the gravitational modulation $\tau_0(R)$
reduces in the laboratory to standard GR time dilation, but $\tau_0\to\infty$ at the collapse
curvature $R_c$ (D4). \textbf{Honesty about scope.} The inflationary sector reproduces,
without improving it, the periodic non-minimal natural inflation of the literature; the
novelty is not there. The rigorous test (E1) shows the clock sector to be non-discriminating
($\varepsilon=1-10^{-112}$): the CMB is the sole judge, and the framework's proper signature
\emph{would be} the cross-sector correlation $\beta(\mathrm{CMB})=\beta(\text{closure})$ ---
a conjecture conditional on the reduction (\S\ref{sec:collapse}), not a theorem. Each
statement is tagged by its actual status.

\smallskip\noindent\textbf{Keywords:} scale duality; scalar-tensor gravity; non-minimal
coupling; natural inflation; regular black holes; intermittent quantum existence; intrinsic
decoherence; relational clock; unification.\\
\textbf{PACS:} 04.50.Kd, 04.70.Bw, 98.80.Cq, 03.65.Ta, 04.62.+v.\quad
\textbf{MSC:} 83C57, 83D05, 83F05, 81P15.

\smallskip\noindent\textit{Epistemic status:} \statut{proof} exact algebra/computer algebra
$\cdot$ \statut{numerical} reproducible integration $\cdot$ \statut{conjecture} plausible,
unproven $\cdot$ \statut{heuristic} analogy $\cdot$ \statut{speculation} open $\cdot$
\statut{open} to be established $\cdot$ \statut{to check} external ref. No conjecture is
presented as a theorem.
\end{abstract}

% =====================================================================
\section{Introduction}
Reconciling quantum mechanics and general relativity has been an open problem for a century.
On the gravitational side, two requirements bear on any Planck-regime physics: removing
singularities [3--7] and accounting for the primordial accelerated expansion [8, 9]. On the
quantum side, decoherence and the selection of a preferred basis must be explained without an
ad hoc external environment [2, 20]. Established programmes (strings, loop gravity,
noncommutative geometry, hybrid semiclassical gravity) address these fronts by distinct
routes.

We explore a single geometric hypothesis. The logarithmic scale axis, usually non-compact, is
taken to be a circle: by stereographic projection, the Planck scale (pole $P$) and a
« unified » scale (pole $\infty^*$) become two poles of a single manifold [1]. There are two
such axes --- spatial ($\psi$) and temporal ($\chi$) --- whence the arena
$\mathcal{M}=M_{3+1}\times S^1_\psi\times S^1_\chi$. The spatial axis governs the
gravitational/cosmological face [1]; the temporal axis, the quantum face --- the time of the
Quantum Existence Discontinuity hypothesis (HDEQ) [2].

\paragraph{An ontological clarification, up front.} HDEQ does not posit a discrete time.
Time remains a continuous parameter; what is intermittent is the active existence of quantum
systems, in cycles of duration $\tau_0$ (active fraction $\varepsilon$). The qualifier
« temporal » for the $\chi$ axis denotes the clock role it plays (\S\ref{sec:quantum}), not a
metric time-like signature.

\paragraph{Scope and epistemic honesty.} This work does not claim to solve unification on its
own. It proposes a coherent and falsifiable direction: modeling decisions are assumed, results
are tagged by their actual status, and it is for the community to test, extend or refute it.

% =====================================================================
\section{Formal framework: the arena and the action}
\begin{postulate}[compactification]
The fundamental manifold is $\mathcal{M}=M_{3+1}\times S^1_\psi\times S^1_\chi$, the scale
axes being circles via $\psi=2\arctan[\ln(\lambda/\ell_P)]$ and
$\chi=2\arctan[\ln(\tau/\tau_P)]$. The poles $\pm\pi$ are identified ($\infty^*$); $0$ is the
Planck point $P$.
\end{postulate}

\subsection{Action, field-space metric (B1) and viable band}
\begin{equation}
S=\int d^4x\,\sqrt{-g}\Big[\tfrac12 fR-\tfrac12\omega\big((\partial\psi)^2+(\partial\chi)^2\big)-V\Big],
\quad f=1+\beta(\cos\psi+\cos\chi),\quad V=\alpha(2-\cos\psi-\cos\chi).
\end{equation}
\statut{proof} (B1) In the Einstein frame, the field-space metric $\mathcal{G}_{ab}$ is
positive-definite $\Leftrightarrow f>0$: a kinetic part $\omega/f$ (positive-definite) plus a
rank-1 positive-semidefinite non-minimal part. The viable band is the open interval
$\beta\in(-0.5;-0.15)$ (at $\beta=-1/2$, $f\to0$ at $P$).
\statut{proof} (F2) $S^1_\psi,S^1_\chi$ are internal circles (Euclidean $\sigma$-model
target), not spacetime dimensions; $\chi$ is not a time-like coordinate $\to$ no ghost, no
closed time-like curves. Winding $=$ the topological scale charge.

% =====================================================================
\section{The gravitational and cosmological sector (« infinitely large »)}\label{sec:cosmo}
\statut{proof} At frozen fields, Einstein gravity emerges; at the symmetric vacuum $P$
($V=0$), the cosmological constant vanishes. At the pole $\infty^*$, the source reduces to a
pure vacuum energy, giving a locally de Sitter space of rate $H^2=2\alpha/[3(1-\beta)]$ (one
axis) or $4\alpha/[3(1-2\beta)]$ (two axes), with finite Kretschmann scalar $K=24H^4$.

In homogeneous cosmology, the canonical field $\varphi=f_s\psi$ sees the potential
$V=\alpha[1-\cos(\varphi/f_s)]$: a natural inflation at locked non-minimal coupling ---
$f=1+\beta\cos\psi$ has the same structure as the potential (studied independently in [16]).
\statut{numerical} Integration starts at $\infty^*$ (de Sitter phase) and rolls down to $P$
($\Lambda=0$), with graceful exit. Slow-roll predictions (Einstein frame):

\begin{table}[h]\centering\small
\begin{tabular}{ccccl}
\toprule
$\beta$ & $f_s$ ($\Mpl$) & $n_s$ & $r$ & status\\
\midrule
$0$ & $7$ & $0.960$ & $0.078$ & $r$ too large\\
$+0.3$ & $7$ & $0.949$ & $0.128$ & excluded\\
$-0.3$ & $5$ & $0.958$ & $0.017$ & compatible (edge)\\
$-0.3$ & $6$ & $0.959$ & $0.027$ & compatible (edge)\\
$-0.3$ & $7$ & $0.960$ & $0.036$ & at the BK18 edge\\
\bottomrule
\end{tabular}
\caption{\small Slow-roll $(n_s,r)$ at $N=55$. Target Planck$+$BK: $n_s=0.9649\pm0.004$,
$r<0.036$.}\label{tab:nsr}
\end{table}

For the natural sign $\beta\ge0$, $r\gtrsim0.08$ is excluded [12, 13]. The viable window is a
narrow band $\beta\in(-0.5;-0.15)$ with \textbf{$f_s\approx5$--$6\,\Mpl$} (open interval, cf.\
II.A), giving $r\approx0.02$--$0.04$ correlated with $n_s\approx0.96$. \statut{numerical}
\emph{Honesty caveat (P1).} Across this band $n_s\approx0.958$--$0.961$, i.e.\ $\sim1$--$2\sigma$
below the Planck central value: the framework lives \emph{at the edge} of the contour. The CMB
data provide a constraint on the scale-duality coupling $\beta$ within this framework (the
$(n_s,r)$ phenomenology being shared with [16, 24, 25]), testable by the Simons Observatory
[27] and LiteBIRD [14] ($\sigma(r)\sim10^{-3}$) by 2030.\footnote{The CMB-S4 project
(DOE--NSF), long the reference instrument for this test, was halted: on 9 July 2025 the DOE
and NSF stated they could no longer support the project (orderly shutdown), which is therefore
no longer funded [15]. The falsifiability stands, carried by the Simons Observatory (Chile,
operating), LiteBIRD ($\sim$2032) and BICEP/Keck (ongoing). \statut{to check}}

\statut{numerical} (A1) Second-order slow roll with the exact field-space metric reproduces
Table~\ref{tab:nsr} to better than $0.5\%$, and fixes $\alpha$ via $A_s=2.1\times10^{-9}$:
$\alpha\approx0.5$--$1.2\times10^{-9}\,\Mpl^4$, an inflationary scale
$V^{1/4}\sim1.1$--$1.4\times10^{16}$~GeV (GUT). At the inflationary $f_s\approx5$--$6$ this
gives $\tau_0=2\pi\sqrt{\omega/\alpha}\approx3.4$--$3.8\times10^{-37}$~s (self-consistent: the
clock period now uses the \emph{same} $f_s$ as inflation, resolving the earlier $f_s=57$
mismatch which gave $r\approx0.14$, excluded). New prediction, stable across the band: the
running $dn_s/d\ln k\approx-6.5\times10^{-4}$ --- compatible with Planck ($-0.004\pm0.007$).
(Remaining: mode-by-mode Mukhanov--Sasaki, expected sub-percent; $(n_s,r,\alpha)$ are computed
on the single-field representative trajectory, the full two-field treatment --- $f$ offset at
$\chi=0$, turning effects --- remaining an untested assumption, track A2.)

\paragraph{Honesty about novelty.} This sector reproduces, without improving it, the periodic
non-minimal natural inflation of Ferreira--Notari--Simeon [16], its scalar-tensor extension
[24], and the $(n_s,r,\text{running})$ analyses against BICEP/Keck [25]. The novelty of the
present work is therefore not the inflationary phenomenology, but the attempted cross-sector
unification (\S\ref{sec:collapse}).

% =====================================================================
\section{The static sector: no regular black hole (closure)}\label{sec:static}
\statut{proof}\statut{numerical} At one field [1], the static radial profile is closed: the
frozen branch $\psi\equiv-\pi$ is exactly Schwarzschild--de Sitter; the rolling branch
develops a singular horizon (collapse of $B=e^{\Phi+\Lambda}$, negative Misner--Sharp mass,
divergent $R$). We rest the closure on the \emph{robust} argument rather than on a
single out-of-band tableau: by Derrick + topology, the static profile (connecting $\infty^*$
to $P$) is an \emph{unprotected half-soliton} $N=\tfrac12$, hence contractible --- no regular
interpolating solution exists, for any $\beta$ in the band. The integer soliton $N=1$ has mass
$8\sqrt{\omega\alpha}\approx2$--$5\times10^{15}$~GeV (GUT scale; $\approx3\times10^{15}$ at the
central $f_s\approx6$, with $\alpha$ fixed in A1) --- a value harmonized with the corrected
$f_s$ (P1).

\paragraph{Un-freezing the temporal axis.} \statut{proof} Making $\chi$ dynamic,
$\chi\equiv0$ remains an exact solution and reproduces the closure, but the coupling of $\chi$
to $R$ renormalizes gravity: on the slice $\chi=0$, $f_{\mathrm{eff}}=(1+\beta)+\beta\cos\psi$,
and the stability threshold shifts ($\beta_c=-1/4$ for two axes). The closure conclusion is
robust; the precise threshold value is reduction-dependent ($-1/3$ mono-axis, $-1/4$ on the
$\chi=0$ slice, $-1/6$ on the diagonal --- see \S\ref{sec:collapse}).

% =====================================================================
\section{The pole $\infty^*$: a structural tension between closure and unwinding}\label{sec:collapse}
\statut{numerical} Does \emph{dynamical} (non-static) collapse cap the curvature? The answer
turns on the nature of the pole $\infty^*$ --- which is \emph{not} reduction-independent. This
is the crux, and we record it as the central structural result of this synthesis: an open
tension, not a settled verdict.

\paragraph{Mono-axis / Jordan reading.} On the single active axis ($\chi$ frozen), for the
Jordan potential $V$, $\infty^*$ is a maximum ($V''(-\pi)=-\alpha<0$). Homogeneous spherical
collapse with the field rolling then drives $\psi$ \emph{through} $\infty^*$: the curvature is
not capped, no regular core forms. In this reading the two verdicts --- graceful exit in
cosmology, no regular core in collapse --- share the cause « $\infty^*$ is a hilltop ».

\paragraph{Two-field / Einstein-frame reading.} But the cosmology that fixes $(n_s,r)$ is,
strictly, a two-field system on the diagonal $\psi=\chi$. \statut{numerical} In the
Einstein-frame diagonal landscape $U=V/f^2$, the nature of $\infty^*$ changes sign at
$\beta=-1/6$: for $\beta>-1/6$ it is a maximum (rolls), but for $\beta<-1/6$ --- i.e.\ across
the \emph{entire viable band} $\beta\in(-0.5;-0.15)$ --- $\infty^*$ is a \emph{trapped local
minimum (a false vacuum)} behind an internal barrier $\theta_b$. In this reading a collapse
driving the field to $\infty^*$ could be \emph{retained} there (de Sitter core) rather than
rolling through --- which would instead point toward a capped curvature, i.e.\ a possible
regular core.

\begin{figure}[h]\centering
\includegraphics[width=0.74\textwidth]{figs_u3_en/fig_pole.png}
\caption{\small \statut{numerical} Einstein-frame diagonal landscape
$U(\theta)/\alpha=2(1-\cos\theta)/(1+2\beta\cos\theta)^2$ ($\psi=\chi=\theta$, exact closed
form). For $\beta<-1/6$ (green, blue --- the entire viable band) $\infty^*$ ($\theta=\pi$) is
a \emph{trapped local minimum} behind a barrier (false vacuum); for $\beta\ge-1/6$ (grey, red)
it is a maximum (the field rolls through). The closure wants $\infty^*$ trapped, the
unwinding wants it traversable: the same pole cannot be both.}\label{fig:pole}
\end{figure}

\paragraph{The tension.} The two readings disagree on the very nature of $\infty^*$, hence on
whether collapse is regularized. They are not cosmetic. The \emph{static} closure rests on
$\infty^*$ being a trap (Derrick: the field cannot connect $\infty^*$ to $P$, so the only
static solution is the frozen Schwarzschild--de Sitter), while the literal « unwinding »
picture (the parchment, the field winding \emph{through} $\infty^*$) requires $\infty^*$ to be
traversable. The same pole cannot be both, and the band $\beta<-1/6$ that delivers the closure
is exactly the band that obstructs dynamical unwinding. We therefore record:
\begin{itemize}
\item the \textbf{static} closure (no regular black hole) is robust --- it follows from
Derrick $+$ topology (contractible half-soliton $N=\tfrac12$) and holds across reductions. \statut{proof}
\item the \textbf{dynamical} collapse verdict is reduction-dependent and \emph{open}:
rolls-through-and-not-capped in the mono-axis/$V$ reading, but a trapped false vacuum (hence a
candidate for a capped, regular core) in the two-field Einstein-frame reading across the
viable band. \statut{conjecture}\statut{open}
\end{itemize}

\paragraph{A failed shortcut.} One might hope that genuine multi-field dynamics (a « spiral »:
the field winding off the diagonal) suppresses $r$ or sustains the unwinding. \statut{numerical}
A full two-field integration with an injected lateral boost shows it does not: the turning is
transient ($\omega/H\approx1.6$ initially) and damped by Hubble friction well before the CMB
window; $r$ stays $\approx0.046$, never suppressed. The reason is structural --- $P$ is a
point well, not a ring valley, so nothing sustains an orbit. The « unwinding » is therefore
best read as a \emph{topological/spatial} statement (the integer soliton charge $N$), and the
literal \emph{temporal} unwinding (a wound $N\ge1$ trajectory through $\infty^*$) as a
\emph{conditional} signature (log-periodic features) available only \emph{outside} the viable
band --- not as an established dynamical mechanism.

\paragraph{The one quantum-sector opening.} If a regularization exists in the viable band, the
candidate is not classical roll-through but the quantum clock: at the collapse curvature
$R_c=2\alpha/|\beta|$, the effective mass of $\chi$ vanishes and $\tau_0\to\infty$ (D4, \S\ref{sec:quantum}) ---
the cycle no longer closes, which could forbid the classical roll-through and retain the field
at $\infty^*$. This is the only avenue by which the unification loop might close. \statut{speculation}

% =====================================================================
\section{The quantum sector and the bridge (« infinitely small »)}\label{sec:quantum}
\subsection{The existence factor exactly derived (D1)}
\statut{proof} The oscillation of $\chi$ in $V$ is a pendulum (sine-Gordon). For an amplitude
$\chi_{\max}$ (elliptic modulus $k=\sin(\chi_{\max}/2)$), the identity for $\langle\mathrm{sn}^2\rangle$
gives the exact form (verified by integrating $\chi''=-\sin\chi$ to $10^{-16}$):
\[
\varepsilon=\langle\cos^2(\chi/2)\rangle=\frac{E(k)}{K(k)},\qquad \tau_0(k)=4\sqrt{\omega/\alpha}\,K(k).
\]
At small amplitude $\varepsilon\to1-\chi_{\max}^2/8$ and $\tau_0\to2\pi\sqrt{\omega/\alpha}$.
Approaching the separatrix ($\infty^*$), $\varepsilon\to0$ (suspension). $\varepsilon$ is
therefore no longer a free parameter: it is $E(k)/K(k)$, fixed by the amplitude.

\begin{figure}[h]\centering
\includegraphics[width=0.66\textwidth]{figs_u3_en/fig_epsilon.png}
\caption{\small \statut{numerical} The existence factor $\varepsilon=E(k)/K(k)$,
$k=\sin(\chi_{\max}/2)$ (verified against the pendulum $\langle\cos^2(\chi/2)\rangle$ to
$10^{-6}$). At small amplitude $\varepsilon\to1-\chi_{\max}^2/8$ (continuous QM); at the
separatrix $\varepsilon\to0$ (suspension). $\varepsilon$ is derived, not free.}\label{fig:eps}
\end{figure}

\subsection{Operatorial realization (D1) and gravitational modulation (D4)}
\statut{conjecture} (realized, D1) $\chi$ is a Page--Wootters relational clock [22]: on
$\mathcal{H}_{\mathrm{clock}}\otimes\mathcal{H}_{\mathrm{sys}}$ with a timeless constraint, the
system evolves with respect to $\chi$ --- defining « temporal » as a clock role, with no
time-like signature (answer to F2). For \emph{constant} $\Hsys$, the evolution gated by
$g(\chi)=\cos^2(\chi/2)$ gives over one period $\hat U(\tau_0)=\exp(-i\Hsys\varepsilon\tau_0)$,
i.e.\ $\Heff=\varepsilon\Hsys$ --- the Floquet result $\Omega_F=\varepsilon\Omega_R$ of [2],
here verified for constant $\Hsys$ (gated qubit, error $\sim10^{-11}$); the non-commuting case
(full Floquet) is estimated negligible but \emph{not} proven. \statut{open} derive the
timeless constraint from the $\chi$ action; exactness of the two-state reduction.

\statut{proof} (D4) The coupling of $\chi$ to $R$ shifts its effective mass:
$\tau_0(R)=2\pi\sqrt{\omega/(\alpha+\tfrac12\beta\Mpl^2 R)}$. Two effects, not to be confused.
(A) $\sqrt{-g_{00}}\approx1+\phi/c^2$: standard gravitational dilation ($\tau_0$ is a
proper-time scalar), observable but not new --- it is GR. (B) the dynamical shift
$\delta\tau_0/\tau_0\approx-(\beta/4)(\Mpl^2 R/\alpha)\propto$ local $R$: it vanishes in vacuum
and is $\sim10^{-81}$ in matter --- unobservable. But at $R_c=2\alpha/|\beta|$ (the collapse
threshold of \S\ref{sec:collapse}), the effective mass of $\chi$ vanishes: $\tau_0\to\infty$,
the quantum clock freezes. The same $R_c$ that triggers the descent makes the cycle diverge ---
a concrete realization of trapping at $\infty^*$ \statut{conjecture}. ($\tau_0(R)$ is
small-amplitude; the exact behavior near $R_c$ requires the full anharmonic potential.)

\begin{figure}[h]\centering
\includegraphics[width=0.66\textwidth]{figs_u3_en/fig_tau0R.png}
\caption{\small \statut{numerical} Gravitational modulation
$\tau_0(R)/\tau_0(0)=1/\sqrt{1-R/R_c}$, $R_c=2\alpha/|\beta|$. The clock period diverges as
the curvature approaches the collapse threshold $R_c$: the quantum clock freezes --- the one
candidate route for the quantum sector to regularize collapse.}\label{fig:tau0R}
\end{figure}

\subsection{Lorentz safety of the quantum time}
\statut{proof} $\tau_0\sim10^{-37}$~s corresponds to $\hbar/\tau_0\approx10^{13}$~GeV; one must
check it induces no Lorentz violation (photon time-of-flight, LHAASO). An intermittent
existence violates Lorentz only through a momentum-dependent dispersion modification. Now (i)
the proper evolution of $\hat H$ is exactly phase-linear to all orders in $\tau_0$
(Proposition 1 of HDEQ); (ii) the Maxwell action is Weyl-invariant in $D=4$, so the photon is
blind to the conformal factor $f(\chi)$. Free-photon dispersion remains exactly
$\omega=|p|c$: no LHAASO signal. The only correction (QED trace anomaly) is logarithmic, not
power-law in $E/E_{\mathrm{QG}}$ --- no linear LIV term. \statut{to check}

\subsection{Cross-sector test (E1): the CMB is the sole judge}
\statut{numerical} On the $\alpha$ fixed in A1, $\tau_0\in[2.5;5.3]\times10^{-37}$~s. Decisive
point: unification removes the freedom in $\varepsilon$. The bridge ties $\varepsilon$ to the
present field displacement, which equals dark energy; numerically $\psi_{\mathrm{today}}\sim10^{-56}$~rad,
so $\varepsilon=1-3\times10^{-112}$. Consequences: the Yb$^+$ effect
$\Delta\Gamma/\Gamma=1-\varepsilon^2\approx10^{-112}$, intrinsic decoherence below the floor,
the resonance at $\Delta E\approx1.2\times10^{13}$~GeV --- the entire HDEQ clock sector is
unobservable. \emph{Caveat (open).} The framework sets $\Lambda=0$ at the vacuum $P$, whereas
observed dark energy is non-zero; the identity « present residual $=$ dark energy » that fixes
$\varepsilon=1-10^{-112}$ is not derived, and the reconciliation with the observed $\Lambda$
remains open (a residual-quintessence reading, $\varepsilon$ slightly below $1$, would carry
$\Lambda>0$ at the cost of fine-tuning). The « clocks non-discriminating » conclusion holds in
both readings. \statut{open}

\paragraph{Verdict E1 (rigorous).} The framework passes the consistency test, but the
confrontation is not symmetric: the clock side is satisfied trivially ($\varepsilon\approx1$).
All the falsifiability burden therefore rests on the CMB ($n_s\approx0.96$,
$r\approx0.02$--$0.04$, running $\approx-6.5\times10^{-4}$). The « same $(\alpha,\omega)$ for
the CMB and the clocks » is met automatically on the clock side --- it should \emph{not} be
presented as a discriminating test.

% =====================================================================
\section{Overview and falsifiability criterion}
The vision is an ansatz, not a result: a single arena whose known laws are its two limiting
regimes. (\emph{Infinitely large:}) slow fields $\to$ general relativity, $\Lambda=0$ at
vacuum, inflation. (\emph{Infinitely small:}) circle dynamics $\chi\to$ intermittent
existence, decoherence (HDEQ). The hinge is the $\chi$ axis, which couples to curvature and
modulates quantum time via gravity. The two faces are joined by a bridge now partially proven
($\varepsilon$ derived, gating included; the full operator construction still open).

The truth criterion is numerical, but asymmetric: the CMB $(n_s,r,\text{running})$ is the
judge; the atomic clocks provide a consistency constraint automatically met ($\varepsilon\approx1$),
not a discriminating test. The signature \emph{proposed} for this framework --- which no
isolated cosmological observable possesses, since $(n_s,r,\text{running})$ are shared with
[16, 24, 25] --- is the cross-sector correlation $\beta(\mathrm{CMB})=\beta(\text{closure})$:
the same coupling should fit inflation and satisfy the black-hole closure thresholds.
\emph{This correlation is a conjecture conditional on the reduction} (\S\ref{sec:collapse}),
not a theorem: with three candidate thresholds ($-1/3,-1/4,-1/6$) « $\beta(\text{closure})$ »
is not yet a single number, so settling the reduction (track A2) is the prerequisite to
turning this into a genuine discriminator. That is where falsifiability would live, testable
via LiteBIRD / Simons Observatory.

% =====================================================================
\section{Epistemic status and limitations}
\begin{table}[h]\centering\small
\begin{tabular}{p{0.66\textwidth}l}
\toprule
Content & Status\\
\midrule
Finite action; GR limit; natural inflation at locked coupling, $(n_s,r)$ and a constraint on
$\beta$; robust \emph{static} closure (Derrick $+$ topology); $\mathcal{G}_{ab}$
positive-definite (B1); $\sigma$-model health, no ghost/CTC (F2); photon Lorentz safety;
$\varepsilon=E(k)/K(k)$ derived; $\alpha$ fixed by $A_s$, running $\approx-6.5\times10^{-4}$
(A1); self-consistent $f_s\approx5$--$6$ and $\tau_0$. & \statut{proof}\statut{numerical}\\
\midrule
$\Heff=\varepsilon\Hsys$ via a Page--Wootters clock, restricted to constant $\Hsys$ (D1);
$\tau_0(R)$ derived, $\tau_0\to\infty$ at $R_c$ (D4); E1: clock sector non-discriminating, CMB
sole judge. & \statut{conjecture} (realized)\\
\midrule
\textbf{Central open tension:} the nature of $\infty^*$ (traversable vs trapped) governs the
\emph{dynamical} collapse verdict and the cross-sector signature; reduction not settled (track
A2). Derive the Page--Wootters constraint from the $\chi$ action; reconcile $\Lambda=0$ vs
observed dark energy; inhomogeneous collapse; quantum trapping at $\infty^*$ (regularization);
Born rule; single-branch selection; mode-by-mode Mukhanov--Sasaki. & \statut{open}\statut{speculation}\\
\bottomrule
\end{tabular}
\end{table}

\noindent The distinction of statuses is not an ornament but a constitutive part of the
method: it ensures this work can be reused without any statement exceeding what is actually
established.

% =====================================================================
\section{Conclusion}
A single scale-duality field, on an arena with compactified axes (internal circles, not
spacetime dimensions), generates a predictive and falsifiable inflationary cosmology, settles
the \emph{static} regular-black-hole question by one structural cause (Derrick $+$ topology),
and offers a bridge toward the quantum time of HDEQ now partially proven: the existence factor
$\varepsilon=E(k)/K(k)$ is derived, the active/inactive mechanism and the Lorentz-safe free
propagation established. The \emph{dynamical} regularization, however, is not settled: it
hinges on the nature of the pole $\infty^*$, traversable in one reduction and a trapped false
vacuum in another --- a tension we record as the central open problem, with the quantum
freezing at $R_c$ as the one candidate route to close the loop. This is not a closure of
unification; it is a direction whose very nature is to be refutable. That is the price --- the
honest tagging of each statement --- at which an individual intuition can become a collective
direction.

\paragraph{Acknowledgements.} The author thanks the symbolic and numerical computation tools
(SymPy, SciPy), which enabled systematic verification at each step, and the
gr-qc/hep-th/quant-ph community for the standards of rigour and falsifiability that guided this
work. The formalization was conducted in dialogue with analytical-assistance systems, under
the author's direction, who bears full responsibility for the scientific content.

% =====================================================================
\section*{Reproducibility}
The numerical results and symbolic verifications are reproduced by two standalone scripts
(NumPy/SymPy/SciPy/Matplotlib), with sections coded so as to be specific to this document
(no collision with other script sets):
\begin{itemize}
\item[{\texttt{[S1]}}] existence factor $\varepsilon=E(k)/K(k)$ verified against the pendulum $\langle\cos^2(\chi/2)\rangle$; small-amplitude $\varepsilon\to1-\chi_{\max}^2/8$; $\tau_0(k)=4\sqrt{\omega/\alpha}\,K(k)$.
\item[{\texttt{[S2]}}] nature of $\infty^*$: diagonal threshold $\beta=-1/6$ (trapped/rolls) and transverse isocurvature mass (sign change at $\beta=-1/4$).
\item[{\texttt{[S3]}}] inflation on the viable band $f_s=5$--$6$: $(n_s,r)$ diagonal vs mono (agreement); $\alpha$, $V^{1/4}$, $\tau_0\approx3.5\times10^{-37}$~s, soliton mass (P1).
\item[{\texttt{[S4]}}] gravitational modulation $\tau_0(R)$; $R_c=2\alpha/|\beta|$; $\delta\tau_0/\tau_0=-(\beta/4)\Mpl^2R/\alpha$.
\item[{\texttt{[S5]}}] spiral test (two-field integration): the turning is damped, $r$ never suppressed.
\end{itemize}
Files: \texttt{U3\_SYNTH\_reproductible.py} (calculations) and \texttt{Synthese\_figures.py}
(Fig.~\ref{fig:eps}, \ref{fig:pole}, \ref{fig:tau0R}). \emph{Convention}: $\Mpl=1$; $\alpha$
cancels in the $(n_s,r)$ ratios. The prefix \texttt{U3\_SYNTH\_} and the \texttt{[S$\cdot$]}
codes are specific to this document.

% =====================================================================
\begin{thebibliography}{99}\small
\bibitem{u1} H.\ Bakkaoui, \textit{Scale-duality field in gravity} (framework, corrections and results), Zenodo (2026).
\bibitem{u2} H.\ Bakkaoui, \textit{Quantum Existence Discontinuity (HDEQ)}, Zenodo (2026).
\bibitem{bardeen} J.\ M.\ Bardeen, Proc.\ GR5 (Tbilisi, 1968), p.\ 174.
\bibitem{hayward} S.\ A.\ Hayward, \textit{Phys.\ Rev.\ Lett.} \textbf{96}, 031103 (2006).
\bibitem{dymnikova} I.\ Dymnikova, \textit{Gen.\ Relativ.\ Gravit.} \textbf{24}, 235 (1992).
\bibitem{penrose} R.\ Penrose, \textit{Phys.\ Rev.\ Lett.} \textbf{14}, 57 (1965).
\bibitem{he} S.\ W.\ Hawking \& G.\ F.\ R.\ Ellis, \textit{The Large Scale Structure of Space-Time} (CUP, 1973).
\bibitem{guth} A.\ H.\ Guth, \textit{Phys.\ Rev.\ D} \textbf{23}, 347 (1981).
\bibitem{ffo} K.\ Freese, J.\ A.\ Frieman \& A.\ V.\ Olinto, \textit{Phys.\ Rev.\ Lett.} \textbf{65}, 3233 (1990).
\bibitem{bek} J.\ D.\ Bekenstein, \textit{Phys.\ Rev.\ D} \textbf{51}, R6608 (1995).
\bibitem{bezr} F.\ L.\ Bezrukov \& M.\ Shaposhnikov, \textit{Phys.\ Lett.\ B} \textbf{659}, 703 (2008).
\bibitem{planck} Planck Collaboration, \textit{Astron.\ Astrophys.} \textbf{641}, A10 (2020).
\bibitem{bk} BICEP/Keck Collaboration, \textit{Phys.\ Rev.\ Lett.} \textbf{127}, 151301 (2021).
\bibitem{litebird} LiteBIRD Collaboration, \textit{Prog.\ Theor.\ Exp.\ Phys.} \textbf{2023}, 042F01 (2023).
\bibitem{cmbs4} K.\ N.\ Abazajian \textit{et al.} (CMB-S4), arXiv:1610.02743 (2016) --- \emph{project halted, July 2025}.
\bibitem{fns} R.\ Z.\ Ferreira, A.\ Notari \& G.\ Simeon, \textit{JCAP} \textbf{11} (2018) 021.
\bibitem{fmm} V.\ P.\ Frolov, M.\ A.\ Markov \& V.\ F.\ Mukhanov, \textit{Phys.\ Rev.\ D} \textbf{41}, 383 (1990).
\bibitem{rv} C.\ Rovelli \& F.\ Vidotto, \textit{Int.\ J.\ Mod.\ Phys.\ D} \textbf{23}, 1442026 (2014).
\bibitem{cr} A.\ Connes \& C.\ Rovelli, \textit{Class.\ Quantum Grav.} \textbf{11}, 2899 (1994).
\bibitem{zurek} W.\ H.\ Zurek, \textit{Rev.\ Mod.\ Phys.} \textbf{75}, 715 (2003).
\bibitem{kaiser} D.\ I.\ Kaiser, \textit{Phys.\ Rev.\ D} \textbf{81}, 084044 (2010).
\bibitem{pw} D.\ N.\ Page \& W.\ K.\ Wootters, \textit{Phys.\ Rev.\ D} \textbf{27}, 2885 (1983).
\bibitem{bars} I.\ Bars, \textit{Class.\ Quantum Grav.} \textbf{18}, 3113 (2001). \statut{to check}
\bibitem{simeon} G.\ Simeon, arXiv:2002.07625 (2020).
\bibitem{bostan} I.\ D.\ Bostan, arXiv:2209.02434 (2022).
\bibitem{mv} C.\ Marletto \& V.\ Vedral, \textit{Phys.\ Rev.\ D} \textbf{95}, 043510 (2017).
\bibitem{so} Simons Observatory Collaboration (P.\ Ade \textit{et al.}), \textit{JCAP} \textbf{02} (2019) 056, arXiv:1808.07445.
\end{thebibliography}

\end{document}
